\chapter{Discussion}
\label{ch:discussion}

This chapter interprets the results of Chapter~\ref{ch:results}, examines what
they do and do not establish, and considers threats to validity, ethical
implications and deployment requirements.

\section{Interpreting the Screening Performance}
\label{sec:disc_ml}

\subsection{What 70.9\% accuracy means here}

Taken in isolation, an accuracy of 70.9\% appears modest against a literature
reporting figures above 98\%. The comparison is not meaningful, for two reasons
established in Chapters~\ref{ch:literature} and \ref{ch:results}. First, those
models consume the laboratory markers from which the label is derived; the
present model is explicitly denied them. Second, they are typically evaluated on
a small, class-balanced benchmark, whereas this evaluation uses an unselected
population in which the high-risk class constitutes 4.5\% of participants.

The more informative framings are these. Against a majority-class baseline of
82.9\% accuracy and 33.3\% balanced accuracy, the model achieves 66.0\% balanced
accuracy and a macro ROC AUC of 0.79---substantial discrimination where a
trivial classifier has none. And the ordinal structure of the errors matters: a
quadratic weighted $\kappa$ of 0.479, with a confusion matrix in which errors are
overwhelmingly adjacent-category, indicates the model has learned a graded risk
surface rather than an arbitrary partition.

\subsection{Was the cascade justified?}

The cascade was adopted to create an explicit control point for the high-to-low
error. On the held-out set, 2 of 38 high-risk participants were assigned to the
low category, and 36 of 38 (94.7\%) received some referral-generating category.
This is the design working as intended: the Stage-1 threshold of 0.44, below the
naive 0.50, widens the gate at a known and accepted cost in precision.

That cost is substantial and should not be understated. Moderate precision is
0.242, meaning roughly three in four moderate predictions are false positives.
A deployment would need to be explicit with users that a moderate or high result
means ``this warrants a confirmatory test'', not ``you probably have kidney
disease''. The system's patient-facing text is written accordingly, but the
underlying statistical reality is a screening instrument tuned for sensitivity,
and its precision must be communicated honestly to both patients and clinicians.

An important caveat on the cascade claim: this thesis reports the cascade's
performance but did not conduct a controlled comparison against a single
three-class classifier trained on the same features and splits. The architectural
argument for the cascade---that it provides a natural place to encode asymmetric
error costs---is sound in principle, but the empirical claim that it outperforms
a well-tuned single model on the high-to-low error specifically remains untested
here, and is identified as future work.

\subsection{The moderate class}

Moderate recall of 0.505 is the weakest metric in the evaluation. The failure is
structurally explicable: moderate is the middle band of a continuous risk
surface, bounded on both sides, and its 52 errors split almost evenly between
the two neighbouring categories. From a clinical-consequence standpoint, the
moderate class is also the least costly place to err---a moderate participant
misclassified as high still receives a referral, and one misclassified as low is
the least severe of the possible missed-referral cases. Nonetheless, this is
where the greatest headroom for improvement lies.

\section{What Actually Delivers Safety}
\label{sec:disc_safety}

The most consequential finding of this work is a negative one about language
models and a positive one about guardrails.

Across all six arms, including both fine-tuned configurations, between 6.3\% and
10.3\% of raw model outputs contained a critical safety violation. Prompting the
model with nine explicit safety rules did not eliminate this. Grounding it in
retrieved clinical guidance did not eliminate this. Fine-tuning it on 540
safety-screened examples did not eliminate this---and the improvement it did
produce ($-3.0$ points) failed to reach significance. Only the deterministic
post-generation validator eliminated it, reducing the delivered rate to zero and
lifting emergency-redirect recall from 65--75\% to 100\%, with both effects
surviving Holm correction at $p<0.001$.

This has a direct architectural implication that generalises beyond kidney
screening. In a patient-facing generative system, safety should be treated as a
property of the \emph{pipeline}, not of the \emph{model}. Techniques that shape
the model's distribution---prompting, retrieval, fine-tuning---improve the prior
but provide no guarantee, because they operate probabilistically. A deterministic
gate applied to the produced text is the only component in this architecture
that can offer a bounded guarantee, and the measurements here indicate it is
carrying essentially the entire safety burden.

It is worth being precise about the scope of that guarantee. The validator
guarantees that no text matching a pattern in its detector set will be delivered.
It guarantees nothing about patterns outside that set. The ten false-positive
defects documented in Section~\ref{sec:corrections} are evidence of this
limitation operating in the opposite direction---rules that were too broad---but
the symmetric risk, of rules that are too narrow, is not directly measurable
from within the same rule system. This is the fundamental epistemic limit of
rule-based safety and is discussed further in Section~\ref{sec:validity}.

\section{The Role of Fine-Tuning}
\label{sec:disc_finetune}

If fine-tuning did not deliver a statistically established safety improvement,
was it worth doing? The evaluation supports a qualified yes, on grounds other
than safety.

The decisive evidence is the format result. The base model given the compact
prompt (arm A) produced correctly structured output in 0\% of cases; the
fine-tuned model given the same prompt (arm F) achieved 100\%. Since the compact
prompt is what makes the 25\% prompt-token reduction possible, fine-tuning is
the enabling condition for that saving. Fine-tuning also reduced destructive
fallbacks ninefold (3.00\% to 0.33\%), meaning the patient receives a
personalised answer rather than a generic template far more often; improved
readability by roughly half a grade level; and---when served with its matched
prompt---raised grounding to 99.7\%.

The honest summary is therefore that fine-tuning bought efficiency, format
reliability and answer quality, and that these are legitimate returns. It did
not buy the safety guarantee, and the thesis does not claim it did. This
distinction is only visible because the ablation was designed to separate the
components; an end-to-end evaluation of the deployed system alone would have
reported 0\% unsafe output and invited the incorrect inference that the
fine-tuned model is safe.

\subsection{The prompt--model matching effect}

An unanticipated and practically useful finding is the grounding collapse when
the fine-tuned model is served with the full prompt: 99.7\% grounding under the
compact prompt it was trained with, against 67--69\% under the full prompt, with
retrieval held identical. Fine-tuning on a specific prompt format appears to
induce a dependence on that format, such that serving the adapter with a
different prompt degrades precisely the behaviours the adapter was trained to
produce. The practical rule---serve a fine-tuned model with the prompt it was
trained against---is obvious in retrospect but is easy to violate in a system
that supports multiple prompt variants, which is why the implementation binds
prompt selection automatically to adapter presence rather than leaving it as an
independent configuration option.

\section{Personalisation}
\label{sec:disc_personal}

Personalisation was pursued through three mechanisms: factor-aware retrieval,
inclusion of the complete patient profile in the prompt, and explicitly computed
per-section personalisation instructions. The measured personalisation score of
1.99 out of 2 across 300 profiles indicates the objective was met, though the
metric is a rule-based proxy for whether the text references the individual's
own reported values, not a human judgement of usefulness.

The development history is instructive. The system initially produced nearly
identical text for every patient in a given risk category, and the cause was
twofold: a safety rule instructing the model to give ``only general wellness
education'', which directly contradicted the personalisation objective, and an
abstract instruction to ``personalise'' that provided nothing concrete to act
on. Rewriting the rule to authorise condition-specific reference while retaining
every prohibition, and computing explicit per-section instructions from the
patient's own values, resolved it. The general lesson is that instructions to a
language model must be operationally specific; ``personalise this'' is not an
instruction a model can reliably execute, whereas ``mention reducing salt because
this person has high blood pressure'' is.

\section{Threats to Validity}
\label{sec:validity}

\subsection{Internal validity}

The chief internal threat is that the safety evaluation and the deployed
validator share a methodological lineage: both are rule-based detectors written
by the same authors. A pattern neither anticipates will be neither blocked nor
counted, which could inflate the apparent safety of the delivered output. The
ten corrected false positives demonstrate that the detectors were imperfect in
the direction of over-triggering; nothing in the design can rule out symmetric
under-triggering. Independent adjudication---ideally clinician review of a
sample of delivered outputs---is required to bound this, and is the single most
valuable next evaluation step.

A second internal threat concerns the label. The target is a KDIGO grid category
computed from measured \eGFR{} and \UACR{} at a single point in time. KDIGO
requires abnormality persisting beyond three months for a CKD diagnosis, which a
cross-sectional survey cannot establish. The label is therefore correctly
described as a laboratory-derived risk stratification rather than a confirmed
diagnosis, and the model's task is to predict that stratification.

\subsection{External validity}

The model is trained on a United States adult survey population. Applying it to
a Bangladeshi population---the setting most relevant to this institution---would
require local validation, because the distribution of the predictors, the
prevalence of the outcome, and the relationships between them may all differ.
Nothing in this work establishes transfer.

The language-model evaluation uses 300 synthetic profiles, scored by the
deployed cascade rather than drawn from real patients. This is appropriate for
measuring the internal behaviour of the safety mechanisms---it allows full
control of the emergency and adversarial subgroups, which could not be
ethically constructed from real patients---but it measures mechanism
effectiveness, not real-world outcome validity. No claim is made that patients
who read this output understand it better or act on it, which would require a
user study.

\subsection{Statistical validity}

The emergency subgroup ($n=60$) and adversarial subgroup ($n=30$) are modest.
A 100\% or 0\% result at these sizes is precise for this benchmark---the Wilson
interval for 60/60 is $[94.0, 100]$---but should not be read as a guarantee at
larger scale. The F vs E comparison in particular has limited power because both
arms sit at the ceiling on the primary metrics.

All reported language-model results use deterministic decoding at a single seed.
Sampling variability across seeds was not characterised for the final
configuration, so the reported figures describe this decoding configuration
rather than an expectation over stochastic decoding.

\subsection{Construct validity}

Two metric definitions deserve scrutiny. \emph{Grounding} measures whether the
model cited at least one identifier that was actually supplied to it. It is a
measure of source presence and citation honesty, not of factual correctness: a
response could cite a genuinely retrieved passage while mischaracterising its
content. \emph{Retrieval relevance} uses knowledge-base topic matching as a
proxy, with no human relevance labelling. Both metrics are therefore weaker than
their names might suggest, and are described accordingly wherever they appear.

\section{Ethical Considerations}
\label{sec:ethics}

\textbf{Risk of false reassurance.} The most serious ethical hazard of a
screening tool is a false negative that discourages a person from seeking care.
The system mitigates this through a sensitivity-first threshold, through
explicit statements that a low result is not a guarantee, and through a
validator rule that blocks reassurance language inconsistent with the computed
risk. It cannot eliminate the hazard.

\textbf{Scope discipline.} The system never diagnoses, never mentions
medication, and never issues individualised nutrient targets. These constraints
are enforced deterministically rather than requested, precisely because the
consequences of violation---a patient stopping an antihypertensive on the advice
of a web application---are severe and irreversible.

\textbf{Data protection.} No real patient data was used in development. NHANES is
public and de-identified, and the fine-tuning corpus is entirely synthetic. At
inference, all computation is local: the language model, embedder and vector
index run on the deployment machine, and no patient input is transmitted to any
external service. This is a direct consequence of choosing a 4B open-weight
model over a hosted frontier API, and it is a meaningful privacy property for
health data.

\textbf{Equity.} The training population is not the deployment population for
this institution's context. Deploying without local validation risks
systematically miscalibrated risk estimates for the served population, which is
an equity concern and not merely a technical one.

\textbf{Automation bias.} A fluent, confident, well-formatted explanation can
induce more trust than its underlying accuracy warrants. The reliability
indicator, the explicit ``screening, not diagnosis'' framing, and the display of
input completeness are all countermeasures, but the tension between
persuasiveness and epistemic honesty is intrinsic to using a language model in
this role.

\section{Deployment Considerations}
\label{sec:disc_deploy}

Three practical gaps stand between the current system and a real deployment.

First, the input schema exposes no free-text symptom field, so the validator's
emergency detection is reachable in production only through the vitals pathway.
The symptom vocabulary is implemented and exercised by the evaluation harness,
but a deployed system intended to catch symptom-reported emergencies requires
that field to be added end to end.

Second, persistence is best-effort by design: a database outage silently skips
the write. This is correct for a research demonstrator, where a failed save must
never block a screening, but a production clinical record system would require
durable queuing.

Third, latency of roughly 31 seconds per full guidance request on the
development hardware is acceptable for a demonstration but not for a
high-throughput clinic. Serving optimisation---batching, quantised inference,
adapter merging---was not pursued, since it does not affect any scientific claim
in this thesis.
